% Reading map prepared for James Avery (2026-10-09): how the GeoVac papers
% build on J.E. Avery & J.S. Avery, "Hyperspherical Harmonics and Their
% Physical Applications" (World Scientific 2017). Content mirrors the
% reviewed web page (claims + citation review, 2026-10-09).
\documentclass[11pt]{article}
\usepackage[T1]{fontenc}
\usepackage[utf8]{inputenc}
\usepackage{amsmath, amssymb}
\usepackage[margin=1.05in]{geometry}
\usepackage{longtable, booktabs, array}
\usepackage{enumitem}
\usepackage{xcolor}
\usepackage{hyperref}
\hypersetup{colorlinks=true, linkcolor=teal!70!black, urlcolor=teal!70!black}
\setlist{itemsep=2pt, topsep=3pt}
\emergencystretch=2em
\newcommand{\tst}[1]{\nolinkurl{#1}}
\makeatletter\g@addto@macro{\UrlBreaks}{\do\_\do\.}\makeatother

\newcommand{\rel}[1]{\textsc{#1}}
\newcommand{\bk}[1]{\textit{#1}}

\title{Labels to Generalized Sturmians:\\ a map of the GeoVac papers onto
\emph{Hyperspherical Harmonics and Their Physical Applications}}
\author{J.~Loutey\thanks{Independent researcher,
\texttt{jloutey@gmail.com}. Developed in an AI-assisted workflow
(implementation and drafting with Anthropic's Claude);\ errors are the
author's. Papers and tests:\
\url{https://github.com/jloutey-hash/geovac}.}}
\date{Prepared for James Avery, October 2026}

\begin{document}
\maketitle

\noindent This note maps what the GeoVac papers build on in
J.~E.~Avery and J.~S.~Avery, \emph{Hyperspherical Harmonics and Their
Physical Applications} (World Scientific, 2017), what they add, and where
your methods meet a quantum computer. It runs from the discrete object
GeoVac starts from, through Fock's projection, up to your
generalized-Sturmian secular equation and its molecular form.

Book references give the chapter and section title, which are the stable
reference;\ equation numbers may differ slightly between printings. GeoVac references give
the paper and the regression test (under \texttt{tests/}) that checks the
claim. Every number is quoted from the papers or their tests. Throughout,
$K$ is the number of configurations (the dimension of the secular matrix),
and the \emph{1-norm} is the sum of coefficient magnitudes when that matrix
is written as a weighted sum of unitaries, which sets the cost of
qubitization.

\paragraph{How each connection is labelled.}
\rel{builds on};\ \rel{extends};\ \rel{puts numbers on} (the book states
it qualitatively and a GeoVac paper measures it);\ \rel{yours, credited}
(results our papers had presented as their own until a direct reading of
your book this month, and which now credit you);\ \rel{independent route};\
\rel{scope of our claim} (reading your work narrowed a GeoVac statement).

%-----------------------------------------------------------------------
\section{The chain, one rung at a time}

\subsection*{1. The discrete object: a label set with a Casimir on it}
Paper~0 starts from an isotropic packing construction. Its shells have
capacities $2k-1$, the same counts as the spherical harmonics $Y_{lm}$, and
grouping them gives $n^2$ states per shell;\ the paper records that
grouping as an observation, not a consequence of its axioms. The object is
the exact label set $(n,l,m)$ of the $S^3$ harmonics, on which the
Laplace--Beltrami operator acts as the diagonal SO(4) Casimir,
\[
-\Delta_{S^3} = L^2 + A^2,\qquad \text{eigenvalue } n^2-1 \text{ on every label},
\]
a classical structure (Pauli 1926, Bargmann 1936) that Paper~7 credits so.
It is \emph{not} the lattice's own graph Laplacian:\ that operator never
changes $l$, splits into $n_{\max}$ disconnected blocks, and does not
converge to $\Delta_{S^3}$ --- Paper~7 proves this and records it as a
negative result.\\
\emph{GeoVac:} Paper~0 abstract;\ Paper~7 \S New contributions;\
\tst{test_paper7_label_set.py}, \tst{test_graph_s3_nonconvergence.py}.\\
\emph{Book} (\rel{independent route}):\ Ch.~2, \bk{Degeneracy of
hyperspherical harmonics}, $(\lambda+1)^2$ for $d=4$, eq.~(2.43);\
\bk{Eigenfunctions of generalized angular momentum}, eq.~(2.8). Paper~0
reaches the count combinatorially.

\subsection*{2. Fock's projection: from labels to the hydrogen spectrum}
Paper~7 audits Fock's stereographic map and the chordal-distance identity
in 18 symbolic checks of the known identities. The $S^3$ Coulomb kernel
has eigenvalue $2\pi^2/n$, which forces $p_0=Z/n$ and gives
$E_n=-Z^2/2n^2$ as a function of the Casimir. The labels carry the
degeneracies and selection rules;\ the projection carries the physics.
GeoVac reads the two as dual descriptions and takes no position on which is
more basic.\\
\emph{GeoVac:} Paper~7 \S Conformal map, \S Energy shell;\
\tst{test_fock_projection.py}, \tst{test_fock_laplacian.py},
\tst{test_paper7_fock_kernel.py}.\\
\emph{Book} (\rel{independent route}):\ Ch.~5, \bk{Fock's original
treatment}, eqs.~(5.10)--(5.29). Paper~7 credits Fock 1935 as the primary
source for the one-electron chain.

\subsection*{3. Coulomb Sturmians at a shared scale}
Fix the scale and let the charge vary instead:\ all functions share one
$k$, so all live on the same $S^3$. Papers~8, 34 and 60 use the
potential-weighted orthonormality and the one-body matrix in its Sturmian
form,
\[
\int \chi^*_{n'l'm'}\,\frac{n}{kr}\,\chi_{nlm}\,d^3r=\delta,\qquad
h = k^2\bigl(1-\tfrac{m}{2}\bigr) - Zk\,\mathrm{diag}(1/n),
\]
with $m$ the overlap. GeoVac's production qubit Hamiltonians use hydrogenic
orbitals with scale $Z/n$;\ those are not a Sturmian basis, and Paper~19
says so.\\
\emph{GeoVac:} Paper~8 \S Sturmian basis;\ Paper~34 \S Sturmian
projection;\ Paper~60 eq.~(pw);\ \tst{test_paper60_sturmian.py}.\\
\emph{Book} (\rel{builds on}):\ Ch.~6, \bk{Potential-weighted
orthonormality}, eqs.~(6.32), (6.39);\ Ch.~9, \bk{Matrices representing
kinetic energy and nuclear attraction}, eqs.~(9.40), (9.44).

\subsection*{4. Many electrons}
For $N$ electrons the same statement holds on $S^{3N-1}$:\ the
Laplace--Beltrami operator is the diagonal SO($3N$) Casimir
$\lambda(\lambda+3N-2)$. Paper~7 cites it as standard hyperspherical-harmonic
theory (J.~Avery, 1989);\ Paper~13 uses the SO(6) case for helium, through
other references.\\
\emph{GeoVac:} Paper~7 \S N-electron;\ Paper~13 \S Casimir;\
\tst{test_algebraic_angular.py}.\\
\emph{Book} (\rel{builds on}):\ Ch.~2, eqs.~(2.1), (2.8);\ the $d=6$
tables in Ch.~3, \bk{Normalization}.

\subsection*{5. Generalized Sturmians:\ the isoenergetic secular equation}
Paper~60 takes your secular equation as written. Fix the energy, solve for
the scale:\ the roots are values of $p_\kappa$, each giving an energy
directly. There is no overlap matrix and no outer energy search:
\[
\bigl[\,T^0 + T' - p_\kappa \mathbb{1}\,\bigr]B = 0,\qquad
T^0 = Z\,\mathcal{R}_\nu\,\delta,\qquad E=-p_\kappa^2/2 .
\]
The form for an arbitrary weight potential $V_0$ is also yours, on
Goscinski's construction. Earlier versions of Paper~60 presented it as
their own;\ it is now credited:
\[
\sum_\nu \langle \Phi_\mu|\,V-\beta_\nu V_0\,|\Phi_\nu\rangle\, C_\nu = 0 .
\]
\emph{GeoVac:} Paper~60 eq.~(secular), eq.~(general $V_0$);\
\tst{test_paper60_sturmian.py}, \tst{test_paper60_general_v0.py}.\\
\emph{Book} (\rel{builds on};\ \rel{yours, credited} for eq.~(7.13)):\
Ch.~7, \bk{Sturmians and the many-particle Schr\"odinger equation},
eq.~(7.13);\ \bk{Derivation of the secular equations}, eqs.~(7.78)--(7.89);\
Goscinski's weighted equation, eq.~(7.1).

\subsection*{6. The tied scale, measured}
In the isoenergetic formulation the basis scale is tied to the eigenvalue.
Your text notes that the resulting Goscinskian exponents come out almost
exactly at their variational optimum, and that the method suits excited
states especially. For the single $1s^2$ configuration the remark is
exact:\ the exponent is the textbook $27/16$ and the energy $-2.847$~Ha.
Paper~60 measures how the departure grows as configurations are added. For
the helium ground state the tied scale stalls at $7.46$~mHa while the same
span with a free scale reaches $1.28$~mHa ($K=130$). At $K=452$ the ground
state sits $6.82$~mHa from exact and the $2\,^1S$ state $1.716$~mHa:\ the
cost is state-dependent and falls on the ground state, as your split-shell
account explains. Paper~60 also proves a variational upper bound for this
formulation, with an inertia corollary, and an interlacing bound;\ we did
not find the upper bound stated in your books or theses.\\
\emph{GeoVac:} Paper~60 eq.~(scale lock), \S ground state;\
\tst{test_paper60_scale_lock.py}, \tst{test_paper60_no_selection.py}.\\
\emph{Book} (\rel{puts numbers on}):\ Ch.~7, closing text of
\bk{Derivation of the secular equations}, and the $^1S$ table for helium.

\subsection*{7. Two centres:\ the Shibuya--Wulfman metric returns}
For atom-centred Sturmians the metric is weighted by the ket's own nucleus,
which makes each centre's block exactly the identity:
\[
S_{\tau'\tau} = \frac{n_\tau}{k}\,\Bigl\langle \chi_{\tau'}\Bigr|\frac{1}{r_{a(\tau)}}\Bigl|\chi_\tau\Bigr\rangle
= \Bigl\langle \chi_{\tau'}\Bigr|\frac{p^2+k^2}{2k^2}\Bigl|\chi_\tau\Bigr\rangle .
\]
Paper~60's regression test reproduces ten of your closed-form
Shibuya--Wulfman integrals to about $10^{-12}$. On a full $\sigma$ basis
($n\le N$, $l<n$, $m=0$ on each centre) the conditioning grows exponentially
with the number of shells:\ $18$ to $1.1\times10^{11}$ from $3$ to $15$
functions per centre at $kR=3$. Paper~60 shows why. In this metric each
centre's set is complete and the other centre's set is its translate, which
is unitary here, so the two-centre metric is a finite section of
\[
\begin{pmatrix} I & U\\ U^* & I\end{pmatrix},\qquad U \text{ unitary},\quad
\text{spectrum } \{0,2\}.
\]
No preconditioner that acts centre by centre can help;\ that is proved. The
remedy you use for the two-centre density expansion, phasing out the
smallest eigenvectors, also works on this metric:\ at threshold $10^{-6}$ it
holds the conditioning at $10^{6}$ with an $\mathrm{H}_2^+$ accuracy loss
below $10^{-7}$~Ha through $36$ functions per centre. The threshold has to
fall as the basis grows.\\
\emph{GeoVac:} Paper~60 \S Molecular frontier;\
\tst{test_paper60_sigma_basis.py}, \tst{test_paper60_sigma_redundancy.py}.\\
\emph{Book} (\rel{puts numbers on}):\ Ch.~8, \bk{Shibuya--Wulfman
orbitals};\ Ch.~9, \bk{Shibuya--Wulfman integrals}, eq.~(9.26), and
\bk{Matrices representing kinetic energy and nuclear attraction} (the
own-nucleus form), eqs.~(9.42)--(9.44);\ \bk{2-center densities}, the remark
after eq.~(9.75) that the density-fitting matrix $1+W$ becomes
ill-conditioned through overcompleteness at small $S$, handled by phasing
out the smallest eigenvectors. The translation reading itself is Shibuya
and Wulfman's (1965).

%-----------------------------------------------------------------------
\section{On a quantum computer:\ generated, loaded, paid for}
The idea we discussed:\ the angular structure can be produced on the device
from quantum-number labels, so only the radial part has to be loaded, and
the radial part depends on one variable. Paper~60 states the idea and tests
what it costs. It does not measure the loading cost itself;\ that remains a
motivating argument.

\paragraph{Generated, from labels.}
\begin{itemize}
\item Angular couplings are exact algebraic numbers (square roots of
rationals):\ Gaunt coefficients, $3j$ and $6j$ symbols, computed from
$(n,l,m)$ with no quadrature.
\item On one centre they are sparse, and the sparsity depends only on
$l_{\max}$, not on the potential (Paper~22). Your Ch.~8 says the same of the
one-centre coefficients:\ large but very sparse.
\item At two centres $l$ stops being conserved and only the axial $m$
survives (Papers~14, 58). Your Ch.~9 notes that the two-centre density
coefficients are not sparse.
\end{itemize}

\paragraph{Loaded, one-dimensional.}
\begin{itemize}
\item The radial factors depend on one variable:\ $r$ for an atom, the
hyperradius for several electrons, and for two centres closed-form
functions of the internuclear distance $R$.
\item If the angular part is generated on the device, the loading cost would
be a small set of radial seeds rather than an $O(M^4)$ integral tensor.
\item For two centres those seeds close over $\{e^aE_1(a),\ \ln,\ \gamma\}$,
classically (Sugiura, Roothaan, Ruedenberg);\ Paper~58 organises them as a
closed seed set. At three centres the $(XY|XZ)$ integral over $1s$ orbitals at
a common exponent becomes an elliptic Bessel moment, genus one, in your
momentum-space setting (Paper~59, which names your 2013 review as the
closest prior art).
\end{itemize}

\paragraph{Paid for:\ the 1-norm.}
\begin{itemize}
\item What sets the quantum cost is chiefly the 1-norm of the encoded
matrix, and for molecules also the metric's conditioning, not the load count
alone.
\item \emph{Atoms:} the secular matrix carries no metric and needs no energy
search. Its 1-norm grows more slowly than its dimension over every basis
computed (helium, $K^{0.82}$ over $K=74$--$164$, local slope falling to
$0.77$ by $K=514$), when the basis grows at fixed $l_{\max}$. It is bought
with the same tied scale that sets the accuracy floor:\ $6.82$~mHa for the
ground state at $K=452$, against $1.716$~mHa for $2\,^1S$ at the same
encoding cost.
\item \emph{Molecules:} the metric returns and the 1-norm is ordinary
polynomial. The metric is paid once, when the one-electron molecular
orbitals are built;\ it does not compound with electron number.
\end{itemize}

\noindent In short:\ your isoenergetic formulation makes the atomic problem
a clean quantum target at a state-dependent accuracy floor, a trade-off
between cost, accuracy and which state is targeted rather than an
unconditional efficiency gain. A targeted search did not find this
combination in the existing quantum-algorithm literature. For molecules, the
cost moves into the Shibuya--Wulfman metric, and rung~7 above is where that
cost now stands.

%-----------------------------------------------------------------------
\section{Concordance:\ your book and the GeoVac papers}
{\small
\renewcommand{\arraystretch}{1.15}
\begin{longtable}{>{\raggedright\arraybackslash}p{0.26\textwidth}>{\raggedright\arraybackslash}p{0.37\textwidth}>{\raggedright\arraybackslash}p{0.12\textwidth}>{\raggedright\arraybackslash}p{0.13\textwidth}}
\toprule
Your book & GeoVac & Relation & Checked by\\
\midrule
\endhead
Ch.~2, generalized angular momentum, eq.~(2.8) & Paper~7, SO($3N$) Casimir (\S N-electron;\ cites J.~Avery 1989) & builds on & \tst{test_algebraic_angular}\\
Ch.~2, degeneracy, $(\lambda+1)^2$, eq.~(2.43) & Paper~0, $n^2$ shell grouping, reached combinatorially & independent & \tst{test_paper7_label_set}\\
Ch.~5, Fock's original treatment, eqs.~(5.10)--(5.18) & Paper~7, stereographic map and chordal identity (credited to Fock 1935) & independent & \tst{test_fock_projection}\\
Ch.~5, kernel eigenvalue and $Z/(kn)=1$, eqs.~(5.21)--(5.29) & Paper~7, kernel eigenvalue $2\pi^2/n$ forces $p_0=Z/n$ & independent & \tst{test_paper7_fock_kernel}\\
Ch.~5, Coulomb Sturmians, eqs.~(5.32), (5.33) & Paper~34, Sturmian reading of the labels at $\lambda=Z/n$ & builds on & \tst{test_paper34_projection_spot_checks_batch1}\\
Ch.~8 and Ch.~9, the four-dimensional three-harmonic coefficient & Papers~32, 38, 44, 45, the Wen--Avery 3-$Y$ integral on $S^3$ & builds on & \tst{test_so4_three_y_integral}\\
Ch.~6, potential-weighted orthonormality, eqs.~(6.32), (6.39) & Paper~60 eq.~(pw);\ Paper~8 Sturmian orthogonality & builds on & \tst{test_paper60_sturmian}\\
Ch.~9, kinetic-energy and nuclear-attraction matrices, eqs.~(9.40), (9.44) & Paper~34, $h=k^2(1-m/2)-Zk\,\mathrm{diag}(1/n)$;\ Paper~60, kinetic matrix $\mathbb{1}-(\text{overlap})/2$ & builds on & \tst{test_paper60_scale_lock} (Paper~60 form)\\
Ch.~7, Goscinski's weighted equation and orthogonality, eqs.~(7.1), (7.5)--(7.8) & Paper~60 \S Isoenergetic reformulation & builds on & \tst{test_paper60_general_v0}\\
Ch.~7, many-particle Schr\"odinger equation, eq.~(7.13) & Paper~60, secular equation for a general weight $V_0$ (now credited to you and Goscinski) & yours, credited & \tst{test_paper60_general_v0}\\
Ch.~7, derivation of the secular equations, eqs.~(7.78)--(7.89) & Paper~60 eq.~(secular) & builds on & \tst{test_paper60_sturmian}\\
Ch.~7, Goscinskian configurations, eqs.~(7.64)--(7.77) & Paper~60:\ He $1s^2=-2.847$~Ha reproduced;\ a domain rule for large $n$ & builds on & \tst{test_paper60_box_rule}\\
Ch.~7, near-optimal exponents;\ suited to excited states & Paper~60:\ ground state $7.46$~mHa tied vs $1.28$~mHa free ($K=130$);\ at $K=452$, $6.82$ vs $1.716$~mHa for $2\,^1S$ & puts numbers on & \tst{test_paper60_scale_lock}\\
Ch.~7, tied-scale formulation & Paper~60:\ variational upper bound (inertia corollary) and interlacing bound;\ upper bound not found in your books or theses & extends & \tst{test_paper60_scale_lock}, \tst{test_paper60_no_selection}\\
Ch.~7, molecular spectra, eqs.~(7.48)--(7.63) & Paper~60, identity-form molecular secular equation;\ $T^0$ diagonal for molecular Sturmians (now credited to you) & yours, credited & ---\\
Ch.~8, Shibuya--Wulfman orbitals:\ fixed-$S$ standard eigenproblem & Paper~60:\ fixed-$S$ vs fixed-$R$ on $\mathrm{H}_2^+$, $8.5\times10^{-2}$ vs $8.5\times10^{-4}$~Ha at 3 functions per centre, same basis & puts numbers on & \tst{test_paper60_sigma_basis}\\
Ch.~9, Shibuya--Wulfman integrals, eqs.~(9.26), (9.27) & Paper~60 eq.~(sw), Hermitian form;\ ten closed forms reproduced to $\sim10^{-12}$ & builds on & \tst{test_paper60_sigma_basis}\\
Ch.~9, nuclear-attraction form, own-nucleus weight, eqs.~(9.42), (9.43) & Paper~60:\ the metric is weighted by the ket's own nucleus (corrects our earlier ``the metric is $V_0$'') & yours, credited & \tst{test_paper60_sigma_redundancy}\\
Ch.~8, Shibuya--Wulfman translation construction & Paper~60:\ two-centre metric a finite section of $\bigl(\begin{smallmatrix}I&U\\U^*&I\end{smallmatrix}\bigr)$, spectrum $\{0,2\}$;\ no centre-wise preconditioner helps (translation reading credited to Shibuya--Wulfman 1965) & extends & \tst{test_paper60_sigma_redundancy}\\
Ch.~9, 2-center densities:\ $1+W$ ill-conditioned by overcompleteness;\ phase out the smallest eigenvectors (text after eq.~(9.75)) & Paper~60:\ the same remedy applied to the Shibuya--Wulfman metric on the $\sigma$ basis;\ at $10^{-6}$ it keeps $\mathrm{H}_2^+$ within $10^{-7}$~Ha through 36 functions per centre & puts numbers on & \tst{test_paper60_sigma_redundancy}\\
Ch.~8 and Ch.~9, one-centre coefficients large but very sparse & Paper~14 Gaunt selection rules;\ Paper~22 sparsity depends only on $l_{\max}$ & independent & Paper~22 / 14 suites\\
Ch.~9, two-centre density coefficients not sparse & Papers~14 and 58:\ $l$ selection lost at two centres, axial $m$ survives & independent & Paper~14 suite\\
Ch.~9, interelectron repulsion between one-centre densities, momentum space, eqs.~(9.52)--(9.60) & Paper~59:\ the $(XY|XZ)$ three-centre integral over $1s$ at a common exponent is an elliptic Bessel moment, genus one (names your 2013 review as closest prior art) & extends & \tst{test_paper59_*}\\
App.~B, Coulomb-Sturmian expansion of functions of $s=kr$, eq.~(B.5);\ STO integrals & Paper~58:\ these integrals are standard Coulomb--Sturmian practice, which GeoVac does not add to & builds on & ---\\
Ch.~8 and Ch.~9, cross-centre overlap couples $n$ & Paper~8:\ our no-binding theorem holds only for a single-$n$ rotation approximation;\ the genuine Sturmian overlap escapes it and binds & scope of our claim & \tst{test_paper8_overlap_cross_n}\\
Ch.~7, generalized Sturmians, eq.~(7.1) & Paper~8 structural theorem, scoped to the approximation, not to your basis & scope of our claim & \tst{test_paper8_sturmian_structural}\\
Ch.~2, hyperradial separation, eq.~(2.1) & Paper~16, kinetic operator in hyperspherical form & builds on & ---\\
Ch.~3, \bk{Normalization}, $d=6$ tables & Paper~13, SO(6) Casimir for helium & independent & \tst{test_algebraic_angular}\\
Ch.~9, Shibuya--Wulfman integrals, own-nucleus form, eq.~(9.43) & Paper~19 \S Shibuya--Wulfman integrals & builds on & ---\\
\bottomrule
\end{longtable}
}

%-----------------------------------------------------------------------
\section{Where it is open}
\begin{description}[style=nextline, leftmargin=1.2em]
\item[The molecular secular equation as a quantum target.] Your
\bk{Molecular spectra} (Ch.~7) gives the identity-form secular equation for
configurations built from molecular generalized Sturmians, with the weights
$\beta_\nu$ found by interpolation. Paper~60 leaves the molecular 1-norm open
with the obstruction identified:\ in the configuration-scaled form the
nuclear-attraction matrix is neither diagonal nor independent of $p_\kappa$
on two centres, and the one measured molecular 1-norm, from an ordinary
two-centre CI, grows polynomially. Whether your molecular-Sturmian form avoids
that obstruction is not tested.
\item[A schedule for phasing out eigenvectors.] Phasing out works on the
$\sigma$ basis. Each fixed threshold tested between $10^{-2}$ and $10^{-4}$
stalls once the basis grows, while $10^{-6}$ holds through 36 functions per
centre. Paper~60 concludes the threshold must fall with the basis;\ a
shrinking schedule has not been tested.
\item[Lowering the ground-state floor.] Paper~60 takes the mechanism from
you:\ in $1s^2$ both electrons share one exponent, so the in-out correlation
that a split-shell $1s\,1s'$ pair carries is missing. What it leaves open is
whether a weighting potential of a different radial shape could lower the
ground-state floor while keeping the Fock mapping.
\item[Three centres.] The three-centre two-electron integral $(XY|XZ)$, over
$1s$ orbitals at a common exponent, is genus one:\ an elliptic Bessel moment
in the momentum-space Slater theory of your Ch.~9 (Paper~59, which names your
2013 review as the closest prior art). Paper~59 identifies the class;\ a
closed form is not known.
\end{description}

%-----------------------------------------------------------------------
\section*{Papers referred to}
All under \texttt{papers/} at \url{https://github.com/jloutey-hash/geovac}.
\begin{itemize}[leftmargin=1.2em]
\item Paper 0, \emph{Angular Momentum as Information Geometry}
\item Paper 7, \emph{The Dimensionless Vacuum}
\item Papers 8--9, \emph{Bond Sphere}
\item Paper 13, \emph{The Hyperspherical Lattice}
\item Paper 14, \emph{Structurally Sparse Qubit Hamiltonians}
\item Paper 16, \emph{The Periodic Table Recast as $S_N$ Representation Theory}
\item Paper 19, \emph{Toward PK-Free Molecular Hamiltonians}
\item Paper 22, \emph{Potential-Independent Angular Sparsity}
\item Papers 32, 38, 44, 45, the spectral-triple papers that use the Wen--Avery integral
\item Paper 34, the projection catalogue
\item Paper 58, \emph{Angular Sparsity Is an Atomic-Sector Property}
\item Paper 59, \emph{The Three-Center ERI Is an Elliptic Bessel Moment}
\item Paper 60, \emph{The Generalized-Sturmian Secular Equation Block-Encodes Without a Metric}
\end{itemize}

\end{document}
